\uuid{y8Bw}
\titre{Limite en l'infini}
\auteur{bodin ; hanani ; mzari}
\organisation{exo7}
\licence{CC BY-NC-SA 4.0}
\datecreate{2020-01-20}
\niveau{L1}
\module{Analyse}
\chapitre{Continuité, limite et étude de fonctions réelles}
\sousChapitre{Limite de fonctions}
\difficulte{1}
\theme{Limites des fonctions réelles, limite en l'infini}
\exo7qcmid{lille-1.limites.005}

\contenu{
\begin{question}
\qclassification{123.03}
\qtags{motcle=limite en l'infini}
Soit $f(x)= \sqrt{x^2+x+1}+x$. Quelles sont les assertions vraies ?
\begin{answers}
    \bad{$\lim_{x\to -\infty} f(x)=+\infty$}

    \bad{$\lim_{x\to -\infty} f(x)=-\infty$}

    \good{$\lim_{x\to -\infty} f(x)=-\frac{1}{2}$}

    \bad{$f$ n'admet pas de limite en $-\infty$.}
\end{answers}
\begin{explanations}
On multiplie le numérateur et le dénominateur de $f$ par l'expression conjuguée de $\sqrt{x^2+x+1}+x$, c'est-à-dire par  $\sqrt{x^2+x+1}-x$. On obtient :
$f(x)= \frac{x+1}{\sqrt {x^2+x+1}-x}$. Attention !  $\sqrt {x^2+x+1}=|x|\sqrt {1+\frac{1}{x}+ \frac{1}{x^2}}= -x\sqrt {1+\frac{1}{x}+ \frac{1}{x^2}}, $ pour $x<0$.
\end{explanations}
\end{question}
}
